\section{Supplementary derivations and theoretical extensions}
\label{sec:appendix}

This appendix collects derivations and complementary results that support the main development
without altering the inferential scope.
The emphasis is on (i) making key steps in the conjugate-exponential-family integration fully
explicit, (ii) clarifying the probabilistic structure of factorizing partition priors, and (iii)
recording a few boundary cases and identifiability remarks that are useful in practice.

\subsection{Conjugate exponential families: detailed block evidence derivation}
\label{app:conjugate}

Section~\ref{sec:ef} states the closed form of the integrated segment evidence
$A_{ij}^{(0)}=p(y_{(i,j]}\mid \text{hyper})$ for exponential-family likelihoods with a conjugate prior,
and gives expressions for posterior moments needed by the Bayes regression curve.
For completeness, we now provide a step-by-step derivation.

\paragraph{Setup.}
Fix a segment (block) $(i,j]$ and write its aggregated sufficient statistics
$T_{ij}=\sum_{t=i+1}^j w_t T(y_t)$ and aggregated exposure $W_{ij}=\sum_{t=i+1}^j w_t$.
Let $\eta\in\mathbb{R}^d$ denote the natural parameter and $A(\eta)$ its cumulant function.
Assume the conditional likelihood factors as
\begin{equation}
 p\big(y_{(i,j]}\mid \eta\big)
 \;=\; \prod_{t=i+1}^j \exp\big\{ w_t\,\langle \eta, T(y_t)\rangle - w_t A(\eta)\big\}\,h(y_t)^{w_t}
 \;=\; \exp\big\{\langle \eta, T_{ij}\rangle - W_{ij}A(\eta)\big\}\,H_{ij},
 \label{eq:app-ef-lik}
\end{equation}
where $H_{ij}:=\prod_{t=i+1}^j h(y_t)^{w_t}$ collects the base-measure terms.
The standard Diaconis--Ylvisaker conjugate prior takes the form
\begin{equation}
 p(\eta\mid \tau,\nu) \;=\; \exp\big\{\langle \eta,\tau\rangle - \nu A(\eta) - \log Z(\tau,\nu)\big\},
 \label{eq:app-ef-prior}
\end{equation}
where $(\tau,\nu)$ are hyperparameters and $Z(\tau,\nu)$ is the normalizer.

\paragraph{Evidence integral.}
Multiplying \eqref{eq:app-ef-lik} and \eqref{eq:app-ef-prior} yields
\begin{align}
 p\big(y_{(i,j]}\mid \eta\big)\,p(\eta\mid\tau,\nu)
 &= H_{ij}\,\exp\Big\{\langle \eta, \tau+T_{ij}\rangle - (\nu+W_{ij})A(\eta) - \log Z(\tau,\nu)\Big\}.
\end{align}
Integrating over $\eta$ therefore gives
\begin{align}
 A^{(0)}_{ij}
 &= \int p\big(y_{(i,j]}\mid \eta\big)\,p(\eta\mid\tau,\nu)\,d\eta \\
 &= H_{ij}\,e^{-\log Z(\tau,\nu)}\int \exp\big\{\langle \eta, \tau+T_{ij}\rangle - (\nu+W_{ij})A(\eta)\big\}\,d\eta \\
 &= H_{ij}\,\frac{Z\big(\tau+T_{ij},\nu+W_{ij}\big)}{Z(\tau,\nu)}.
 \label{eq:app-evidence}
\end{align}
This is the same expression as in Theorem~\ref{thm:ef-integral}, written at the level of a fixed block.

\paragraph{Posterior distribution and moments.}
Bayes' rule shows that the posterior remains in the conjugate family:
\begin{equation}
 p(\eta\mid y_{(i,j]},\tau,\nu) \;\propto\; \exp\big\{\langle \eta,\tau+T_{ij}\rangle - (\nu+W_{ij})A(\eta)\big\}
 \;=\; p\big(\eta\mid \tau+T_{ij},\nu+W_{ij}\big).
 \label{eq:app-post}
\end{equation}
All posterior moments required by BayesBreak can be obtained by differentiating
$\log Z(\tau,\nu)$.
In particular, when $A(\eta)$ is the cumulant generating function,
$\nabla_{\tau}\log Z(\tau,\nu)$ yields the prior mean of the expectation parameter,
and higher derivatives yield higher cumulants.
The block-specific moment numerators $A^{(1)}_{ij},A^{(2)}_{ij}$ used in
Section~\ref{sec:ef} are obtained by multiplying these posterior moments by the evidence $A^{(0)}_{ij}$.

\subsection{Factorizing partition priors as renewal processes}
\label{app:renewal}

Assumption~\ref{ass:factorizing-prior} states that, conditional on the segment count $k$,
the boundary vector $t=(t_0,\dots,t_k)$ has a density proportional to a product of segment-length factors.
This class coincides with renewal-process models on the discrete index set, and admits a hidden-Markov
interpretation that helps motivate the DP recursion.

\begin{definition}[Renewal-process segmentation prior]
\label{def:renewal}
Fix $k\ge1$.
Let $L_1,\dots,L_k$ be positive integer-valued segment lengths with unnormalized mass function
$\varphi(\ell)\propto g\big(x_{i+\ell}-x_i\big)$ for all admissible starts $i$.
A \emph{renewal segmentation} draws lengths $(L_q)_{q=1}^k$ and sets boundaries
$t_q := \sum_{r=1}^q L_r$ with the constraint $t_k=n$.
\end{definition}

\begin{proposition}[Equivalence of renewal and factorized boundary priors]
\label{prop:renewal-equivalence}
Under Definition~\ref{def:renewal}, the induced distribution on $t=(t_0,\dots,t_k)$ satisfies
$p(t\mid k)\propto \prod_{q=1}^k g(x_{t_q}-x_{t_{q-1}})$.
Conversely, any prior of the form \eqref{eq:lengthprior} can be realized as a renewal segmentation prior
with appropriate unnormalized length mass $\varphi$.
\end{proposition}

\begin{proof}
For a fixed boundary vector $t$, the corresponding length sequence is uniquely determined as
$L_q=t_q-t_{q-1}$.
The renewal construction assigns to that sequence mass proportional to
$\prod_{q=1}^k \varphi(L_q)$.
By definition of $\varphi$, this is proportional to $\prod_q g(x_{t_q}-x_{t_{q-1}})$, which proves the first claim.
For the converse, take $\varphi(\ell):=g(x_{i+\ell}-x_i)$ (for any admissible start $i$) as an unnormalized length mass;
then the renewal prior reproduces \eqref{eq:lengthprior} up to the shared normalizer $C_k$.
\end{proof}

\begin{remark}[Hidden-Markov interpretation]
A renewal segmentation can be represented as an HMM whose latent state tracks the ``age'' within the current segment
(or, equivalently, the last boundary location).
The forward recursion in Section~\ref{sec:dp} (cf.\ Theorem~\ref{thm:dp-correctness}, Steps~1--2) coincides with the standard forward algorithm after integrating out
segment parameters into block evidences.
We do not rely on this viewpoint for correctness, but it provides useful intuition for numerical implementations
and for extensions (e.g., adding covariate-dependent hazards).
\end{remark}

\subsection{Posterior boundary and segment-membership marginals}
\label{app:marginals}

Theorem~\ref{thm:dp-correctness} gives closed forms for boundary marginals and for the segment-membership probability
$p\big(t\in(i,j],\,\text{segment }(i,j]\mid y,k\big)$.
We spell out the algebra that leads to these expressions.

\begin{proposition}[Boundary marginals from prefix/suffix evidences]
\label{prop:boundary-marginals}
Fix $k\ge1$ and suppose Assumption~\ref{ass:factorizing-prior} holds.
Let $\widetilde{L}_{r,j}$ and $\widetilde{R}_{r,j}$ be the prefix and suffix evidences defined in \eqref{eq:LR}.
Then for any $t\in\{1,\dots,n-1\}$,
\begin{equation}
 p(t\text{ is a boundary}\mid y,k)
 \;=\; \sum_{r=1}^{k-1}\ \frac{\widetilde{L}_{r,t}\,\widetilde{R}_{k-r,t}}{\widetilde{L}_{k,n}}.
 \label{eq:app-boundary-marginal}
\end{equation}
\end{proposition}

\begin{proof}
By definition, $t$ is a boundary at position $r$ iff $t_r=t$.
For fixed $r$, summing the joint posterior over all boundary vectors with $t_r=t$ yields
\begin{align}
 p(t_r=t\mid y,k)
 &\propto \sum_{\substack{0<t_1<\cdots<t_{k-1}<n\\ t_r=t}}
 \prod_{q=1}^k \widetilde{A}^{(0)}_{t_{q-1},t_q}.
\end{align}
The constraint $t_r=t$ splits the product into a prefix part (segments $1{:}r$ ending at $t$) and a suffix part (segments $r+1{:}k$ starting at $t$).
Summing over all admissible prefixes gives $\widetilde{L}_{r,t}$ by definition of the DP; summing over all admissible suffixes gives $\widetilde{R}_{k-r,t}$.
Normalizing by the total evidence $\widetilde{L}_{k,n}$ yields \eqref{eq:app-boundary-marginal}.
Finally, $p(t\text{ is a boundary}\mid y,k)=\sum_{r=1}^{k-1}p(t_r=t\mid y,k)$, giving the result.
\end{proof}

\subsection{Hazard prior boundary cases}
\label{app:hazard-cases}

The length-aware prior family includes hazard-process priors as a special case.
It is sometimes useful to understand the limiting behavior as the hazard vanishes or becomes certain.

\begin{proposition}[Degenerate hazard limits]
\label{prop:hazard-limits}
Consider the index-uniform design $x_t=t$ and a constant hazard $\rho\in(0,1)$ that induces a geometric length prior
$g(\ell)=\rho(1-\rho)^{\ell-1}$ on segment lengths (up to the constraint that the last segment ends at $n$).
Then:
\begin{enumerate}[label=(\roman*)]
\item As $\rho\downarrow 0$, the prior concentrates on $k=1$ (a single segment) and $p(t\mid k)$ becomes irrelevant.
\item As $\rho\uparrow 1$, the prior concentrates on $k=n$ with unit-length segments, i.e., $t_q=q$ for all $q$.
\end{enumerate}
\end{proposition}

\begin{proof}
For the geometric prior, $\mathbb{P}(L=1)=\rho$ and $\mathbb{P}(L>1)=(1-\rho)$.
As $\rho\downarrow 0$, $\mathbb{P}(L>1)\to 1$ so the probability that a new segment starts at any interior index vanishes,
which forces a single segment spanning $\{1,\dots,n\}$.
As $\rho\uparrow 1$, $\mathbb{P}(L=1)\to 1$ and therefore all segments have length one almost surely, giving $t_q=q$.
\end{proof}

\subsection{Latent groups: identifiability and label switching}
\label{app:label-switching}

In the latent-group extension (Section~\ref{sec:latent-em}), the group labels are arbitrary.
When the prior on segmentations and within-group hyperparameters is exchangeable across groups,
posterior inference is invariant under permutations of the group labels.
This is the familiar \emph{label switching} phenomenon.

\begin{proposition}[Permutation invariance]
\label{prop:perm-invariance}
Assume a $G$-group latent model with an exchangeable prior over group-specific parameters and segmentations:
$p\big((t^{(g)},\psi_g)_{g=1}^G\big)=p\big((t^{(\sigma(g))},\psi_{\sigma(g)})_{g=1}^G\big)$ for every permutation $\sigma$.
Then the joint posterior over $(t^{(g)},\psi_g)_{g=1}^G$ and group assignments $z_{1:S}$ satisfies the same invariance.
In particular, any MAP assignment is non-unique up to label permutation.
\end{proposition}

\begin{proof}
The likelihood factors by subjects given assignments and group-specific parameters, and the prior is invariant by assumption.
Therefore the unnormalized posterior is invariant under simultaneous permutation of group labels in both parameters and assignments.
Normalizing preserves invariance.
\end{proof}

\begin{remark}
Permutation invariance does not prevent accurate segmentation or clustering; it only implies that the label names carry no meaning.
For reporting, one may impose an ordering constraint (e.g., by evidence, segment count, or a canonical parameter summary) or
post-process posterior samples to align labels.
The EM algorithm in Section~\ref{sec:latent-em} can converge to any one of the symmetric modes.
\end{remark}
